\chapter{Рекомендації щодо модернізації}

\section{Ключові напрями оновлення}

На основі проведеного аналізу пропонуються такі напрями модернізації
освітніх програм з дизайну:

\begin{enumerate}
    \item \textbf{Інтеграція AI-інструментів} --- впровадження генеративних
    AI-технологій у навчальний процес
    \item \textbf{Посилення UX/UI компоненти} --- збільшення годин на
    проектування інтерфейсів та прототипування
    \item \textbf{Візуалізація даних} --- додавання окремого модуля з
    інфографіки та дашбордів
    \item \textbf{Етика AI} --- критичне осмислення використання AI у дизайні
\end{enumerate}

\section{Пропозиція навчального плану}


\begin{table}[H]
\centering
\caption{Пропонована структура навчального плану (бакалавр)}
\label{tab:curriculum_proposal}
\begin{tabular}{lccl}
\toprule
Рік & Кредитів & AI-компоненти & Фокус \\
\midrule
Рік 1 & 60 & 4 & Foundations: Traditional skills + Digita... \\
Рік 2 & 60 & 7 & Core Skills: Digital tools mastery + UX/... \\
Рік 3 & 60 & 6 & Advanced: AI integration + Specializatio... \\
Рік 4 & 60 & 3 & Synthesis: Thesis project + Professional... \\

\bottomrule
\end{tabular}
\end{table}


\section{Впровадження}

Рекомендації щодо впровадження змін:
\begin{itemize}
    \item Поетапне оновлення протягом 2-3 років
    \item Підвищення кваліфікації викладачів з AI-інструментів
    \item Партнерство з індустрією для актуалізації змісту
    \item Регулярний перегляд програм (щорічно)
\end{itemize}
